\documentclass[11pt,a4paper]{article}
% ==============================================================
%  Préambule commun aux fiches -- Advanced Probabilities (M1)
%  Usage, dans chaque fiche :
%     \documentclass[11pt,a4paper]{article}
%     % ==============================================================
%  Préambule commun aux fiches -- Advanced Probabilities (M1)
%  Usage, dans chaque fiche :
%     \documentclass[11pt,a4paper]{article}
%     % ==============================================================
%  Préambule commun aux fiches -- Advanced Probabilities (M1)
%  Usage, dans chaque fiche :
%     \documentclass[11pt,a4paper]{article}
%     \input{../preambule.tex}
%     \fiche{C5}{Vecteurs gaussiens}{Gaussian vectors}
% ==============================================================
\usepackage[T1]{fontenc}
\usepackage[utf8]{inputenc}
\usepackage[french]{babel}
\usepackage{lmodern}
\usepackage[a4paper,margin=1.9cm,top=1.7cm,bottom=1.9cm]{geometry}
\usepackage{amsmath,amssymb,amsthm,mathtools}
\usepackage[inline]{enumitem}
\usepackage{booktabs,array,tabularx}
\usepackage{xcolor}
\usepackage[most]{tcolorbox}
\usepackage{fancyhdr}
\usepackage{hyperref}
\hypersetup{colorlinks=true,linkcolor=blue!50!black,urlcolor=blue!50!black}

\setlist{itemsep=1pt,topsep=3pt,parsep=0pt}
\setlength{\parindent}{0pt}
\setlength{\parskip}{3pt}
\emergencystretch=3em

% ---------- Identification de la fiche -------------------------
\newcommand{\ficheid}{}
\newcommand{\fichetitre}{}
\newcommand{\fiche}[3]{%
  \renewcommand{\ficheid}{#1}\renewcommand{\fichetitre}{#2}%
  \pagestyle{fancy}\fancyhf{}%
  \fancyhead[L]{\small\textsf{Advanced Probabilities -- Fiche #1}}%
  \fancyhead[R]{\small\textsf{#2}}%
  \fancyfoot[C]{\small\thepage}%
  \begin{center}
    {\LARGE\bfseries Fiche #1 -- #2}\par\smallskip
    {\large\itshape (#3)}\par\smallskip
    {\small Cours de V.\,M.~Lambert, M1 -- semestre 1}
  \end{center}\vspace{-2pt}\hrule\vspace{6pt}}

% ---------- Environnements numérotés (un seul compteur) --------
\theoremstyle{definition}
\newtheorem{definition}{Définition}
\renewcommand{\thedefinition}{\ficheid.\arabic{definition}}
\newtheorem{theoreme}[definition]{Théorème}
\newtheorem{proposition}[definition]{Proposition}
\newtheorem{lemme}[definition]{Lemme}
\newtheorem{corollaire}[definition]{Corollaire}
\newtheorem{exemple}[definition]{Exemple}
\theoremstyle{remark}
\newtheorem{remarque}[definition]{Remarque}

% Résultat hors cours (standard, signalé comme tel)
\theoremstyle{definition}
\newtheorem{horscours}[definition]{Hors cours}

% Idée de preuve
\newenvironment{idee}{\par\small\textit{Idée de preuve.}~}{\par}

% ---------- Encadrés -----------------------------------------
\tcbset{fichebox/.style={enhanced,breakable,boxrule=0.6pt,arc=2pt,
  left=5pt,right=5pt,top=3pt,bottom=3pt,fonttitle=\bfseries}}
\newtcolorbox{piege}[1][]{fichebox,colback=red!4,colframe=red!60!black,
  title={Pièges classiques},#1}
\newtcolorbox{liens}[1][]{fichebox,colback=blue!3,colframe=blue!50!black,
  title={Liens entre notions},#1}
\newtcolorbox{methode}[1][Méthode]{fichebox,colback=green!4,colframe=green!40!black,title={#1}}
\newtcolorbox{alerte}[1][]{fichebox,colback=orange!6,colframe=orange!70!black,
  title={Attention à l'énoncé},#1}

% Exercice de TD (fiches de rappels)
\newcounter{exo}
\newenvironment{exo}[1][]{\refstepcounter{exo}\par\medskip
  \noindent{\large\bfseries Exercice~\theexo}\ifx\relax#1\relax\else~{\large\bfseries -- #1}\fi\par\nopagebreak\smallskip}{\par}
\newcommand{\outils}{\textbf{Outils du cours.}~}
\newcommand{\meth}{\textbf{Méthode / astuce.}~}
\newcommand{\indic}{\textbf{Indication.}~}
\newcommand{\rep}{\textbf{Réponse.}~}

% ---------- Notations du cours ---------------------------------
\newcommand{\R}{\mathbb{R}}
\newcommand{\N}{\mathbb{N}}
\newcommand{\Q}{\mathbb{Q}}
\newcommand{\C}{\mathbb{C}}
\newcommand{\PP}{\mathbb{P}}
\newcommand{\QQ}{\mathbb{Q}}
\newcommand{\E}{\mathbb{E}}
\newcommand{\F}{\mathcal{F}}
\newcommand{\G}{\mathcal{G}}
\newcommand{\Hc}{\mathcal{H}}
\newcommand{\Tc}{\mathcal{T}}
\newcommand{\Bc}{\mathcal{B}}
\newcommand{\Cc}{\mathcal{C}}
\newcommand{\Ec}{\mathcal{E}}
\newcommand{\Pc}{\mathcal{P}}
\newcommand{\Nc}{\mathcal{N}}
\newcommand{\LN}{\mathcal{LN}}
\newcommand{\Expo}{\mathcal{E}}
\newcommand{\Poi}{\mathcal{P}}
\newcommand{\Ff}{\mathbb{F}}
\newcommand{\1}{\mathbf{1}}
\newcommand{\dd}{\,\mathrm{d}}
\DeclareMathOperator{\Var}{\mathbb{V}}
\DeclareMathOperator{\Cov}{Cov}
\DeclareMathOperator{\card}{card}
\DeclareMathOperator{\diag}{diag}
\newcommand{\indep}{\perp\!\!\!\perp}
\newcommand{\ps}{\ \PP\text{-p.s.}}
\newcommand{\tendloi}{\xrightarrow[n\to+\infty]{\text{loi}}}
\newcommand{\tendps}{\xrightarrow[n\to+\infty]{\text{p.s.}}}
\newcommand{\tendproba}{\xrightarrow[n\to+\infty]{\PP}}
\newcommand{\tendLp}{\xrightarrow[n\to+\infty]{L^p}}
\newcommand{\tend}{\xrightarrow[n\to+\infty]{}}

%     \fiche{C5}{Vecteurs gaussiens}{Gaussian vectors}
% ==============================================================
\usepackage[T1]{fontenc}
\usepackage[utf8]{inputenc}
\usepackage[french]{babel}
\usepackage{lmodern}
\usepackage[a4paper,margin=1.9cm,top=1.7cm,bottom=1.9cm]{geometry}
\usepackage{amsmath,amssymb,amsthm,mathtools}
\usepackage[inline]{enumitem}
\usepackage{booktabs,array,tabularx}
\usepackage{xcolor}
\usepackage[most]{tcolorbox}
\usepackage{fancyhdr}
\usepackage{hyperref}
\hypersetup{colorlinks=true,linkcolor=blue!50!black,urlcolor=blue!50!black}

\setlist{itemsep=1pt,topsep=3pt,parsep=0pt}
\setlength{\parindent}{0pt}
\setlength{\parskip}{3pt}
\emergencystretch=3em

% ---------- Identification de la fiche -------------------------
\newcommand{\ficheid}{}
\newcommand{\fichetitre}{}
\newcommand{\fiche}[3]{%
  \renewcommand{\ficheid}{#1}\renewcommand{\fichetitre}{#2}%
  \pagestyle{fancy}\fancyhf{}%
  \fancyhead[L]{\small\textsf{Advanced Probabilities -- Fiche #1}}%
  \fancyhead[R]{\small\textsf{#2}}%
  \fancyfoot[C]{\small\thepage}%
  \begin{center}
    {\LARGE\bfseries Fiche #1 -- #2}\par\smallskip
    {\large\itshape (#3)}\par\smallskip
    {\small Cours de V.\,M.~Lambert, M1 -- semestre 1}
  \end{center}\vspace{-2pt}\hrule\vspace{6pt}}

% ---------- Environnements numérotés (un seul compteur) --------
\theoremstyle{definition}
\newtheorem{definition}{Définition}
\renewcommand{\thedefinition}{\ficheid.\arabic{definition}}
\newtheorem{theoreme}[definition]{Théorème}
\newtheorem{proposition}[definition]{Proposition}
\newtheorem{lemme}[definition]{Lemme}
\newtheorem{corollaire}[definition]{Corollaire}
\newtheorem{exemple}[definition]{Exemple}
\theoremstyle{remark}
\newtheorem{remarque}[definition]{Remarque}

% Résultat hors cours (standard, signalé comme tel)
\theoremstyle{definition}
\newtheorem{horscours}[definition]{Hors cours}

% Idée de preuve
\newenvironment{idee}{\par\small\textit{Idée de preuve.}~}{\par}

% ---------- Encadrés -----------------------------------------
\tcbset{fichebox/.style={enhanced,breakable,boxrule=0.6pt,arc=2pt,
  left=5pt,right=5pt,top=3pt,bottom=3pt,fonttitle=\bfseries}}
\newtcolorbox{piege}[1][]{fichebox,colback=red!4,colframe=red!60!black,
  title={Pièges classiques},#1}
\newtcolorbox{liens}[1][]{fichebox,colback=blue!3,colframe=blue!50!black,
  title={Liens entre notions},#1}
\newtcolorbox{methode}[1][Méthode]{fichebox,colback=green!4,colframe=green!40!black,title={#1}}
\newtcolorbox{alerte}[1][]{fichebox,colback=orange!6,colframe=orange!70!black,
  title={Attention à l'énoncé},#1}

% Exercice de TD (fiches de rappels)
\newcounter{exo}
\newenvironment{exo}[1][]{\refstepcounter{exo}\par\medskip
  \noindent{\large\bfseries Exercice~\theexo}\ifx\relax#1\relax\else~{\large\bfseries -- #1}\fi\par\nopagebreak\smallskip}{\par}
\newcommand{\outils}{\textbf{Outils du cours.}~}
\newcommand{\meth}{\textbf{Méthode / astuce.}~}
\newcommand{\indic}{\textbf{Indication.}~}
\newcommand{\rep}{\textbf{Réponse.}~}

% ---------- Notations du cours ---------------------------------
\newcommand{\R}{\mathbb{R}}
\newcommand{\N}{\mathbb{N}}
\newcommand{\Q}{\mathbb{Q}}
\newcommand{\C}{\mathbb{C}}
\newcommand{\PP}{\mathbb{P}}
\newcommand{\QQ}{\mathbb{Q}}
\newcommand{\E}{\mathbb{E}}
\newcommand{\F}{\mathcal{F}}
\newcommand{\G}{\mathcal{G}}
\newcommand{\Hc}{\mathcal{H}}
\newcommand{\Tc}{\mathcal{T}}
\newcommand{\Bc}{\mathcal{B}}
\newcommand{\Cc}{\mathcal{C}}
\newcommand{\Ec}{\mathcal{E}}
\newcommand{\Pc}{\mathcal{P}}
\newcommand{\Nc}{\mathcal{N}}
\newcommand{\LN}{\mathcal{LN}}
\newcommand{\Expo}{\mathcal{E}}
\newcommand{\Poi}{\mathcal{P}}
\newcommand{\Ff}{\mathbb{F}}
\newcommand{\1}{\mathbf{1}}
\newcommand{\dd}{\,\mathrm{d}}
\DeclareMathOperator{\Var}{\mathbb{V}}
\DeclareMathOperator{\Cov}{Cov}
\DeclareMathOperator{\card}{card}
\DeclareMathOperator{\diag}{diag}
\newcommand{\indep}{\perp\!\!\!\perp}
\newcommand{\ps}{\ \PP\text{-p.s.}}
\newcommand{\tendloi}{\xrightarrow[n\to+\infty]{\text{loi}}}
\newcommand{\tendps}{\xrightarrow[n\to+\infty]{\text{p.s.}}}
\newcommand{\tendproba}{\xrightarrow[n\to+\infty]{\PP}}
\newcommand{\tendLp}{\xrightarrow[n\to+\infty]{L^p}}
\newcommand{\tend}{\xrightarrow[n\to+\infty]{}}

%     \fiche{C5}{Vecteurs gaussiens}{Gaussian vectors}
% ==============================================================
\usepackage[T1]{fontenc}
\usepackage[utf8]{inputenc}
\usepackage[french]{babel}
\usepackage{lmodern}
\usepackage[a4paper,margin=1.9cm,top=1.7cm,bottom=1.9cm]{geometry}
\usepackage{amsmath,amssymb,amsthm,mathtools}
\usepackage[inline]{enumitem}
\usepackage{booktabs,array,tabularx}
\usepackage{xcolor}
\usepackage[most]{tcolorbox}
\usepackage{fancyhdr}
\usepackage{hyperref}
\hypersetup{colorlinks=true,linkcolor=blue!50!black,urlcolor=blue!50!black}

\setlist{itemsep=1pt,topsep=3pt,parsep=0pt}
\setlength{\parindent}{0pt}
\setlength{\parskip}{3pt}
\emergencystretch=3em

% ---------- Identification de la fiche -------------------------
\newcommand{\ficheid}{}
\newcommand{\fichetitre}{}
\newcommand{\fiche}[3]{%
  \renewcommand{\ficheid}{#1}\renewcommand{\fichetitre}{#2}%
  \pagestyle{fancy}\fancyhf{}%
  \fancyhead[L]{\small\textsf{Advanced Probabilities -- Fiche #1}}%
  \fancyhead[R]{\small\textsf{#2}}%
  \fancyfoot[C]{\small\thepage}%
  \begin{center}
    {\LARGE\bfseries Fiche #1 -- #2}\par\smallskip
    {\large\itshape (#3)}\par\smallskip
    {\small Cours de V.\,M.~Lambert, M1 -- semestre 1}
  \end{center}\vspace{-2pt}\hrule\vspace{6pt}}

% ---------- Environnements numérotés (un seul compteur) --------
\theoremstyle{definition}
\newtheorem{definition}{Définition}
\renewcommand{\thedefinition}{\ficheid.\arabic{definition}}
\newtheorem{theoreme}[definition]{Théorème}
\newtheorem{proposition}[definition]{Proposition}
\newtheorem{lemme}[definition]{Lemme}
\newtheorem{corollaire}[definition]{Corollaire}
\newtheorem{exemple}[definition]{Exemple}
\theoremstyle{remark}
\newtheorem{remarque}[definition]{Remarque}

% Résultat hors cours (standard, signalé comme tel)
\theoremstyle{definition}
\newtheorem{horscours}[definition]{Hors cours}

% Idée de preuve
\newenvironment{idee}{\par\small\textit{Idée de preuve.}~}{\par}

% ---------- Encadrés -----------------------------------------
\tcbset{fichebox/.style={enhanced,breakable,boxrule=0.6pt,arc=2pt,
  left=5pt,right=5pt,top=3pt,bottom=3pt,fonttitle=\bfseries}}
\newtcolorbox{piege}[1][]{fichebox,colback=red!4,colframe=red!60!black,
  title={Pièges classiques},#1}
\newtcolorbox{liens}[1][]{fichebox,colback=blue!3,colframe=blue!50!black,
  title={Liens entre notions},#1}
\newtcolorbox{methode}[1][Méthode]{fichebox,colback=green!4,colframe=green!40!black,title={#1}}
\newtcolorbox{alerte}[1][]{fichebox,colback=orange!6,colframe=orange!70!black,
  title={Attention à l'énoncé},#1}

% Exercice de TD (fiches de rappels)
\newcounter{exo}
\newenvironment{exo}[1][]{\refstepcounter{exo}\par\medskip
  \noindent{\large\bfseries Exercice~\theexo}\ifx\relax#1\relax\else~{\large\bfseries -- #1}\fi\par\nopagebreak\smallskip}{\par}
\newcommand{\outils}{\textbf{Outils du cours.}~}
\newcommand{\meth}{\textbf{Méthode / astuce.}~}
\newcommand{\indic}{\textbf{Indication.}~}
\newcommand{\rep}{\textbf{Réponse.}~}

% ---------- Notations du cours ---------------------------------
\newcommand{\R}{\mathbb{R}}
\newcommand{\N}{\mathbb{N}}
\newcommand{\Q}{\mathbb{Q}}
\newcommand{\C}{\mathbb{C}}
\newcommand{\PP}{\mathbb{P}}
\newcommand{\QQ}{\mathbb{Q}}
\newcommand{\E}{\mathbb{E}}
\newcommand{\F}{\mathcal{F}}
\newcommand{\G}{\mathcal{G}}
\newcommand{\Hc}{\mathcal{H}}
\newcommand{\Tc}{\mathcal{T}}
\newcommand{\Bc}{\mathcal{B}}
\newcommand{\Cc}{\mathcal{C}}
\newcommand{\Ec}{\mathcal{E}}
\newcommand{\Pc}{\mathcal{P}}
\newcommand{\Nc}{\mathcal{N}}
\newcommand{\LN}{\mathcal{LN}}
\newcommand{\Expo}{\mathcal{E}}
\newcommand{\Poi}{\mathcal{P}}
\newcommand{\Ff}{\mathbb{F}}
\newcommand{\1}{\mathbf{1}}
\newcommand{\dd}{\,\mathrm{d}}
\DeclareMathOperator{\Var}{\mathbb{V}}
\DeclareMathOperator{\Cov}{Cov}
\DeclareMathOperator{\card}{card}
\DeclareMathOperator{\diag}{diag}
\newcommand{\indep}{\perp\!\!\!\perp}
\newcommand{\ps}{\ \PP\text{-p.s.}}
\newcommand{\tendloi}{\xrightarrow[n\to+\infty]{\text{loi}}}
\newcommand{\tendps}{\xrightarrow[n\to+\infty]{\text{p.s.}}}
\newcommand{\tendproba}{\xrightarrow[n\to+\infty]{\PP}}
\newcommand{\tendLp}{\xrightarrow[n\to+\infty]{L^p}}
\newcommand{\tend}{\xrightarrow[n\to+\infty]{}}

\begin{document}
\fiche{TD3}{Indépendance, convergence et vecteurs gaussiens}{Independence, Convergence and Gaussian Vectors}

\textit{À lire avant de chercher le TD. Renvois : « C4.5 » = fiche C4, énoncé n°5. Fiches associées : \textbf{C3}, \textbf{C4}, \textbf{C5}.}

\begin{exo}[Borel--Cantelli et Bernoulli]
\outils $\limsup/\liminf$ (C3.5) ; Borel--Cantelli 1 et 2 (C3.6).

\meth Comme $X_n\in\{0,1\}$, « $X_n\to0$ » signifie « $X_n=0$ à partir d'un certain rang » : exprimer l'événement avec $A_n=\{X_n=1\}$.

\indic $\{X_n\to0\}=\liminf_n\{X_n=0\}=\bigl(\limsup_n A_n\bigr)^c$ avec $\PP(A_n)=p_n$. Sens $\Leftarrow$ : BC1. Sens $\Rightarrow$ : contraposée
et BC2, en utilisant que les $A_n$ sont indépendants (fonctions des $X_n$, C3.10).
\end{exo}

\begin{exo}[Somme de gaussiennes indépendantes]
\outils Fonction caractéristique de $\Nc(m,\sigma^2)$ (C5.1) ; $\varphi_{X+Y}=\varphi_X\varphi_Y$ si $X\indep Y$ (C3.15) ; $\varphi$ caractérise la loi (C4.2).

\meth Produit des fonctions caractéristiques (c'est la preuve de C5.2), ou récurrence sur $N$ avec le regroupement par paquets (C3.11) :
$\sum_{n=1}^{N}X_n\indep X_{N+1}$.

\rep $\sum_{n=1}^NX_n\sim\Nc\bigl(\sum m_n,\ \sum\sigma_n^2\bigr)$ (les \emph{variances} s'additionnent).
\end{exo}

\begin{exo}[Limite de gaussiennes]
\outils Théorème de Lévy (C4.5) ; Lemme C4.6.

\meth Écrire $\varphi_{X_n}(t)=e^{itm_n-t^2\sigma_n^2/2}$ et passer à la limite pour $t$ fixé.

\rep $X_n\tendloi\Nc(m,\sigma^2)$ ; si $\sigma=0$, la limite est $\delta_m$ (gaussienne dégénérée).
\end{exo}

\begin{exo}[Théorème de Poisson, binomiale $\to$ Poisson]
\outils Loi binomiale (C2.11) ; critère de convergence en loi pour des v.a.\ à valeurs dans $\N$ (C4.7, hors cours) ; ou bien Lévy (C4.5)
avec l'exemple C4.12.

\meth Écrire $\PP(X_n=k)=\binom nkp_n^k(1-p_n)^{n-k}$ et séparer les trois facteurs, à $k$ fixé.

\indic $\binom nk\sim\frac{n^k}{k!}$, donc $\binom nkp_n^k\to\frac{\lambda^k}{k!}$ ; $(1-p_n)^{n}=\exp(n\ln(1-p_n))$ avec $\ln(1-p_n)\sim-p_n$
($p_n\to0$) ; et $(1-p_n)^{-k}\to1$.
\end{exo}

\begin{exo}[Poisson renormalisée $\to$ gaussienne]
\outils $\varphi$ de $\Poi(\lambda)$ : $e^{\lambda(e^{it}-1)}$ (C4.2) ; $\varphi_{aX+b}$ (C4.2) ; Lévy (C4.5).

\meth Calculer $\varphi_{(X_\lambda-\lambda)/\sqrt\lambda}(t)=\exp\bigl(\lambda(e^{it/\sqrt\lambda}-1)-it\sqrt\lambda\bigr)$ puis développer
$e^{iu}=1+iu-\frac{u^2}2+o(u^2)$ avec $u=t/\sqrt\lambda\to0$.

\indic Les termes en $it\sqrt\lambda$ se compensent ; il reste $-\frac{t^2}2+o(1)$. (La convergence pour $\lambda\to\infty$ réel se ramène
à des suites $\lambda_n\to\infty$.)
\end{exo}

\begin{exo}[Retraits à la banque : Tchebychev contre TCL]
\outils Bienaymé--Tchebychev (C4.11, hors cours) ; TCL (C4.9) ; variance d'une somme indépendante (C3.15) ; exemple C4.13.

\meth $S=\sum_{i=1}^{200}X_i$, $\E S=200\times40$, $\Var S=200\times10^2$. On cherche $c$ tel que $\PP(S\le c)\ge0{,}95$.
Tchebychev : $\PP(S>\E S+a)\le\PP(|S-\E S|\ge a)\le\frac{\Var S}{a^2}$, imposer $\le0{,}05$.
TCL : $\frac{S-\E S}{\sqrt{\Var S}}\approx\Nc(0,1)$ et $c=\E S+1{,}645\sqrt{\Var S}$.

\rep $\E S=8000$, $\sqrt{\Var S}=100\sqrt2\approx141{,}4$.
Tchebychev : $a=\sqrt{20000/0{,}05}\approx632{,}5$, soit $c\approx8633$~€.
TCL : $c\approx8000+1{,}645\times141{,}4\approx8233$~€ (approximation ; bien moins coûteux).
\end{exo}

\begin{exo}[$Z=XY$ : gaussienne sans vecteur gaussien]
\outils Déf.\ de vecteur gaussien (C5.3) ; indépendance (C3.8) ; covariance et $\E(XY)$ (C3.15) ; lien « Liens entre notions » de C5.

\meth
\begin{enumerate}
  \item Conditionner sur la valeur de $Y$ et utiliser la symétrie $-X\sim X$ : $\PP(Z\le z)=p\PP(X\le z)+(1-p)\PP(-X\le z)$.
  \item Tester une combinaison linéaire : $X+Z=X(1+Y)$ vaut $0$ sur $\{Y=-1\}$. A-t-elle un atome sans être constante ?
  \item Comparer $|Z|$ et $|X|$ : trouver $A,B$ avec $\PP(X\in A,Z\in B)\neq\PP(X\in A)\PP(Z\in B)$.
  \item $\Cov(X,Z)=\E(X^2Y)=\E(X^2)\E(Y)$ par indépendance.
\end{enumerate}

\rep 1. $Z\sim\Nc(0,1)$ pour tout $p$. 2. Pour $0<p<1$, non ($X+Z$ a un atome en 0 sans être constante) ; pour $p\in\{0,1\}$,
$Z=\mp X$ et $(X,Z)$ est gaussien (dégénéré). 3. Non (ex.\ $\PP(|X|\le1,|Z|>1)=0$). 4. $\Cov(X,Z)=2p-1=0\iff p=\frac12$ :
\emph{non corrélées mais dépendantes}.
\end{exo}

\begin{exo}[Indépendance de deux combinaisons linéaires]
\outils Transformation affine d'un vecteur gaussien (C5.5) ; Thm des blocs (C5.12).

\meth $(a^tX,b^tX)=AX$ avec $A$ de lignes $a^t,b^t$ : c'est un vecteur gaussien, donc indépendance $\iff$ covariance nulle.

\rep $a^tX\indep b^tX\iff\sum_{i=1}^na_ib_i=0$ (vecteurs $a$ et $b$ orthogonaux).
\end{exo}

\begin{exo}[$N$ indépendant de $Y$]
\outils C5.5, C5.12, $a^tX\sim\Nc(a^tm,a^t\Sigma a)$ (C5.6).

\meth $(N,Y_1,Y_2,Y_3)=AX$ est un vecteur gaussien ; calculer les trois covariances $\Cov(N,Y_j)$ par bilinéarité
($\Cov(X_i,X_j)=\1_{i=j}$).

\rep $N\sim\Nc(0,3)$ ; $\Cov(N,Y_j)=0$ pour $j=1,2,3$, donc $N\indep Y$.
\end{exo}

\begin{exo}[Modèle binomial et formule de Black--Scholes]
\begin{alerte}
\begin{itemize}
  \item L'énoncé \emph{définit} « erf » comme la fonction de répartition de $\Nc(0,1)$ (notée $\Phi$ en général) : notation non
  standard, à ne pas confondre avec la fonction erreur usuelle $\operatorname{erf}_{\text{us}}(x)=\frac2{\sqrt\pi}\int_0^xe^{-t^2}\dd t$
  (lien : $\Phi(x)=\frac12\bigl(1+\operatorname{erf}_{\text{us}}(x/\sqrt2)\bigr)$). Elle vérifie $\operatorname{erf}(x)+\operatorname{erf}(-x)=1$.
  \item IV.2 : il manque un facteur $K$ ; la limite correcte est $Ke^{-rT}\operatorname{erf}(-d_2)-s_0\operatorname{erf}(-d_1)$.
  De même, en V.2, la liste « $r,T,s_0,d_1,d_2$ et erf » omet $K$.
  % Vérifié sur le PDF de l'énoncé (p. 5) : IV.2 écrit e^{-rT} erf(-d2) - s0 erf(-d1), sans K.
\end{itemize}
\end{alerte}
\outils Espace produit et probabilité produit (C3.12) ; indépendance et espérance d'un produit (C3.14, C3.15) ; fonctions caractéristiques
et Lévy (C4.2, C4.5) ; $(1+\frac{z_m}m)^m\to e^z$ (C4.7) ; définition de la convergence en loi avec $f$ continue bornée (C4.3) ; transfert (C2.8).

\meth Modèle produit $\Rightarrow$ $Y_m=\ln(S^{(m)}_T/s_0)$ est une somme i.i.d. ; développement limité de $\Phi_m$ à l'ordre $\frac1m$,
puis Lévy. Le prix limite est $\E(f(Y))$ avec $f$ continue bornée, calculé par complétion du carré gaussien.

\indic
\begin{itemize}
  \item \textbf{I.} Une période : $\Lambda^{(m)}=\{b_m,h_m\}$, $\Tc^{(m)}=\Pc(\Lambda^{(m)})$, $\QQ^{(m)}(\{h_m\})=p_m$.
  $m$ périodes : $\Omega^{(m)}=(\Lambda^{(m)})^m$, $\F^{(m)}=\Pc(\Omega^{(m)})$, $\PP^{(m)}=(\QQ^{(m)})^{\otimes m}$ ; les coordonnées $U^{(m)}_k$ sont
  \textbf{i.i.d.\ de loi $\QQ^{(m)}$}.
  \item \textbf{II.1} Récurrence. \textbf{II.2} Par indépendance, $\E^{(m)}\bigl(\prod_k(1+U_k)\bigr)=(1+\E U_1)^m$ ; comme les deux nombres
  $1+\E U_1$ et $1+\frac{rT}m$ sont $>0$, on obtient $\E U_1=\frac{rT}m$, i.e.\ $p_m\sigma^+-q_m\sigma^-=0$.
  \item \textbf{III.1} $\ln(1+x)=x-\frac{x^2}2+O(x^3)$ avec $x=h_m$ ou $b_m=O(m^{-1/2})$.
  \textbf{III.2} $\Phi_m(u)=p_me^{iu\ln(1+h_m)}+q_me^{iu\ln(1+b_m)}$ et $e^{ix}=1+ix-\frac{x^2}2+o(x^2)$ ; le terme en $\frac1{\sqrt m}$ disparaît grâce à
  $p_m\sigma^+=q_m\sigma^-$, et $p_m(\sigma^+)^2+q_m(\sigma^-)^2=\sigma^+\sigma^-$.
  \textbf{III.3} $Y_m=\sum_k\ln(1+U_k)$, somme i.i.d. \textbf{III.4} $(1+\frac{z_m}m)^m\to e^z$, puis Lévy.
  \item \textbf{IV.1} $x\mapsto(K-s_0e^x)_+$ est continue et \emph{bornée par $K$} : on applique directement la définition de la convergence en loi ;
  et $(1+\frac{rT}m)^{-m}\to e^{-rT}$. Écrire la limite $Y=\mu+\vartheta U$ avec $U\sim\Nc(0,1)$.
  \item \textbf{IV.2} L'intégrande est non nulle ssi $u<-d_2$ ; pour le terme en $s_0$, compléter le carré : $\mu+\vartheta^2/2=rT$ et
  $d_1=d_2+\vartheta$.
  \item \textbf{V.1} $(x-K)_+-(K-x)_+=x-K$ ; linéarité ; $\E^{(m)}(S^{(m)}_T)=s_0(1+\frac{rT}m)^m$ (II.2).
  \textbf{V.2} Parité + IV.2 + $\operatorname{erf}(x)+\operatorname{erf}(-x)=1$.
\end{itemize}

\rep II.2 : $p_m=\frac{\sigma^-}{\sigma^++\sigma^-}$. III.3 : $\varphi_{Y_m}=\Phi_m^m$. III.4 : $Y_m\xrightarrow[m\to+\infty]{\text{loi}}\Nc(\mu,\vartheta^2)$.
V.2 : $\lim C_0^{(m)}=s_0\operatorname{erf}(d_1)-Ke^{-rT}\operatorname{erf}(d_2)$ (formule de Black--Scholes).
\end{exo}

\section*{Synthèse : type de question $\to$ méthode}
\renewcommand{\arraystretch}{1.3}
\begin{tabularx}{\linewidth}{@{}>{\raggedright\arraybackslash}p{0.36\linewidth}X@{}}
\toprule
Type de question & Méthode\\\midrule
Probabilité qu'une infinité de $A_n$ se produisent & Écrire l'événement en $\limsup/\liminf$ ; BC1 (série convergente), BC2 (indépendance) (C3.6).\\
Loi d'une somme de v.a.\ indépendantes & Produit des fonctions caractéristiques (ou génératrices) puis identification (C3.15, C4.2).\\
Convergence en loi & Lévy : limite simple de $\varphi_{X_n}(t)$ (C4.5) ; à valeurs entières : $\PP(X_n=k)\to\PP(X=k)$ (C4.7).\\
Renormalisation $\frac{X-\E X}{\sigma}$ & Développement limité de $\varphi$ à l'ordre 2 ; $(1+\frac{z_n}n)^n\to e^z$.\\
Dimensionner un seuil à 95\,\% & Tchebychev (borne sûre mais grossière) ou TCL (approximation, quantile $1{,}645$) (C4.13).\\
Un couple est-il un vecteur gaussien ? & Oui : l'écrire $AX+b$ avec $X$ gaussien (C5.5). Non : exhiber une combinaison linéaire non gaussienne.\\
Indépendance dans un vecteur gaussien & Covariances croisées nulles (C5.12), \emph{après} avoir justifié le caractère gaussien du vecteur joint.\\
Covariance nulle mais dépendance & Contre-exemple hors vecteur gaussien (ex.~7).\\
Limite d'un prix $\E(f(Y_m))$ & $f$ continue bornée + convergence en loi (C4.3) ; puis calcul gaussien par complétion du carré.\\\bottomrule
\end{tabularx}

\end{document}
